\chapter{洁净室、污染与工艺纪律}

\begin{chaptergoal}
理解洁净室不是“更干净的房间”，而是为了控制颗粒、有机物、金属交叉污染、湿度/静电和流程可重复性；建立 KID 加工所需的记录习惯。
\end{chaptergoal}

\section{先区分四类敌人}

对 KID 来说，“脏”并不是一个单一概念。至少应区分：颗粒导致断线或短路；有机残留改变界面与刻蚀；金属交叉污染改变材料性质；水汽和表面氧化改变界面状态。不同问题需要不同证据，不能把所有低 $Q_i$ 都笼统归因于“加工不干净”。

\processchain{污染/颗粒 $\rightarrow$ 薄膜或界面缺陷 $\rightarrow$ 局部电流拥挤/额外损耗/几何缺陷 $\rightarrow$ $Q_i$、yield、noise 异常}

\section{工艺纪律比单个 recipe 更重要}

第一次进入平台时，优先学会：样品如何编号、不同材料的工具是否分区、哪些设备允许哪些材料、样品如何暂存、何时允许用镊子/真空笔、怎样记录 recipe 版本、异常如何报备。重复性来自流程控制，而不是“记住上次旋钮调到多少”。

\begin{measurebox}
建议从第一批样品就建立 traveler：每一步都写下输入片号、输出片号、设备、recipe、开始/结束时间、异常、是否返工以及下一步允许条件。
\end{measurebox}
